\documentclass[10pt]{article}
%\documentstyle[10pt]{article}
\pagestyle{empty}
\topmargin=0mm
\oddsidemargin=0mm
\textwidth=170mm
\textheight=250mm
\headheight=0mm
\headsep=0mm

\begin{document}
\large
\begin{center}
{\Large \bf
Application of computer algebra systems to construct
high\,-\,order difference schemes
}\\
\vspace {4mm}
A. V. Shapeev
\\
{\it \small Novosibirsk State University, 630090, Novosibirsk, Russia\\
 e-mail: shapeew@itam.nsc.ru}
\end{center}
\vspace {6mm}

 In paper presented, in contrast to \cite{bk:val,bk:shap3}, new approach
to application of computer algebra systems (CAS) to constructing
high\,-\,order difference schemes is stated.
 Programming language is developed, which is some extension of C++ language
and which preserves syntax of MATHEMATICA language.
 Let us call it C\&MATH.
 Language proposed allows to use functions of MATHEMATICA language
(for instance, functions of solving systems of linear algebraic equations,
 simplification of algebraic expressions, etc.).
 Computer program\,-\,preprocessor is written in CAS MATHEMATICA which 
 converts code written in C\&MATH to code in C++ language.
 Then the code obtained can be compiled as ordinary computer program in
 C++ language.

Two ways to construct schemes by method of undetermined coefficients
are implemented in C\&MATH language:
as analytical formulae and numerically \cite{bk:shap1,bk:shap2},
when for given equation and boundary conditions for each computational node
numerical values of difference equations coefficients are found during
solving boundary\,-\,value problem.
 In the first way of scheme constructing, program\,-\,preprocessor forms
 system of linear algebraic equations to determine scheme coefficients on
given stencils and then solves it.
 In the second way, only coefficients of the latter system of equations are
computed explicitly. Further they are used to find scheme coefficients during
solving difference problem.
 Approach proposed was successfully used to get over huge analytical work
concerned with constructing high\,-\,order schemes for solving 
boundary\,-\,value problems for elliptic equations in domain with curvilinear
bounary.
 For the first time we succeeded in creation schemes up to 18th order on
grids with square and triangle cells inside the domain and with irregular
cells near its curvilinear boundary.
 Application of this approach to Navier\,---\,Stokes equations permits to
create 4th\,-\,order scheme on 5\,-\,nodes stencil.
 Great number of numerical experiments on solving different boundary\,-\,value
problems by schemes created and comparison of results obtained with tests and
with known results of other researches were made.
 Effectiveness of obtained schemes was shown \cite{bk:shap1,bk:shap2}.

{\footnotesize This work was supported by FSP "Integration", grant No~276
%{\footnotesize ������ ���������� ������� ��� "����������" \No 276.}
and by the Russian Foundation for Basic
Research, grant No~00-15-96162.}

\begin{thebibliography}{}
\bibitem{bk:val}
 A.N.~Valiullin, V.P.~Shapeev et al.
 Problem of automatic development and investigation of difference schemes 
 in analytic form using computer //
 Dokl. AN SSSR, 1984, Vol.~275, No~3, p.~528--532.

\bibitem{bk:shap3}
V.P. Shapeev.
Application of computer algebra in numerical methods //
Proceedings of 31-th regional youth conference "Problems of theoretical
and applied mechanics" --- Ekaterinburg, 2000. --- Pp.~80--84
(in Russian)

\bibitem{bk:shap1}
 A.V.~Shapeev, V.P.~Shapeev.
 High\,-\,order difference schemes for solving elliptic equations in domain
with curvilinear boundary //
 J. Comput. Math. and Math. Phys., 2000, Vol.~40, No~2 (in Russian).

\bibitem{bk:shap2}
A.V.~Shapeev, V.P.~Shapeev.
High-precision approximation schemes in a domain with curvilinear
boundary. // ENUMATH'99,
Book of abstracts, Jyvaskyla, Finland, July. 26-30, 1999, pp.~134--135.
\end{thebibliography}

\end{document}
